\documentclass[9pt]{beamer}
\usepackage[UTF8]{ctex}

\usepackage{xcolor}
\usepackage{layouts}
\usepackage{minted}
\usepackage[most]{tcolorbox}
\tcbuselibrary{minted}
\usepackage{booktabs}
\usepackage{outlines}
\usepackage{hyperref}
\usepackage{graphicx}
\beamertemplatenavigationsymbolsempty
\usepackage{fontspec}
\setmonofont{Fira Code}[Contextuals=Alternate]


\AtBeginSection[]{
  \begin{frame}
  \tableofcontents[currentsection, hideothersubsections]
  \end{frame}
}

% color
\definecolor{coreFontMain}{RGB}{26,26,26}
\definecolor{coreFontSub}{RGB}{106,106,106}
\definecolor{coreBackgroundLight}{RGB}{230,230,230}
\definecolor{coreBackground}{RGB}{255,255,255}

% heading
\newlength{\headerheight}
\setlength{\headerheight}{0.094\paperheight}


\setbeamerfont{frametitle}{size=\fontsize{12}{14.4}}
\setbeamertemplate{frametitle}
{%
\nointerlineskip
\begin{beamercolorbox}[wd=\paperwidth,left,sep=0pt]{}
\begin{tcolorbox}[enhanced,frame hidden, 
    borderline south={0.5pt}{0pt}{coreFontSub},
    colback=coreBackground, coltext=coreFontSub,height=\headerheight,notitle,left*=0.05\paperwidth,boxsep=0pt,bottom=0pt,halign=left,valign=bottom,before skip=0pt,after skip=0pt, before upper=\strut]
    \usebeamercolor{frametitle}\insertframetitle
\end{tcolorbox}
\end{beamercolorbox}
}


% footer
\newlength{\footerheight}
\setlength{\footerheight}{0.04\paperheight}
\setbeamercolor{footer}{bg=coreBackgroundLight, fg=coreFontMain}

\setbeamertemplate{footline}
{%
\begin{beamercolorbox}[sep=0pt,wd=\paperwidth,center]{footer}
\begin{minipage}[c][\footerheight]{\textwidth}
    \centering
    \insertframenumber/\inserttotalframenumber
\end{minipage}
\end{beamercolorbox}
}



% Title Page

% header
\newlength{\titleheaderheight}
\setlength{\titleheaderheight}{0.0666667\paperheight}
\setbeamercolor{header}{bg=coreBackgroundLight}

%% logo
\newlength{\logoheight}
\setlength{\logoheight}{0.297\paperheight}

%% title
\setbeamercolor{title}{fg=coreFontMain}
\setbeamerfont{title}{size=\fontsize{18.14}{21.768}}
\newlength{\titleheight}
\setlength{\titleheight}{0.19\paperheight}


%% meta
\newlength{\metaheight}
\setlength{\metaheight}{0.21\paperheight}
\setbeamercolor{meta}{fg=coreFontSub}
\setbeamerfont{meta}{size=\fontsize{9.07}{10.884}}

%% contact
\newlength{\contactheight}
\setlength{\contactheight}{0.193333\paperheight}



\setbeamertemplate{title page}{
\nointerlineskip
\begin{beamercolorbox}[wd=\paperwidth,ht=\titleheaderheight]{header}
\end{beamercolorbox}
\nointerlineskip
\begin{beamercolorbox}[wd=\paperwidth,ht=\logoheight]{}
\end{beamercolorbox}
\nointerlineskip
\begin{tcolorbox}[enhanced,frame hidden, borderline south={0.5pt}{0pt}{coreFontMain}, colback=coreBackground,height=\titleheight,boxsep=0pt,bottom=20pt,notitle,halign=flush center,valign=center,before skip=0pt,after skip=0pt, coltext=coreFontMain]
\usebeamercolor{title}
\usebeamerfont{title}\inserttitle

\end{tcolorbox}

\usebeamercolor[fg]{meta}{
\begin{minipage}[c][\metaheight]{\textwidth}
    \centering
    \insertauthor \linebreak
    \insertinstitute \linebreak
    \insertdate
\end{minipage}
}


\begin{tcolorbox}[enhanced,frame hidden, colback=coreBackground,height=\contactheight,notitle,halign=flush center,valign=top,before skip=0pt,after skip=0pt,top=-5pt]
\begin{tabular}{ccc}
\includegraphics[height=3mm]{bohanbeamerstyle/figures/mail.png} & \includegraphics[height=3mm]{bohanbeamerstyle/figures/github.png} &\includegraphics[height=3mm]{bohanbeamerstyle/figures/Twitter.png}
\end{tabular}

\end{tcolorbox}




\begin{beamercolorbox}[sep=2pt,center,wd=\paperwidth,ht=\footerheight]{footer}
\tiny
\insertframenumber/\inserttotalframenumber
\end{beamercolorbox}
\nointerlineskip
}


% 
\setbeamercolor{normal text}{bg=coreBackground,fg=coreFontMain}

% itemize
\setbeamertemplate{itemize items}[circle]

\setlength{\leftmarginii}{10pt}
\setbeamercolor{itemize item}{fg=coreFontMain}
\setbeamercolor{itemize subitem}{fg=coreFontMain}
\setbeamercolor{itemize subsubitem}{fg=coreFontMain}



% code

\newtcblisting{codebox}[2]{listing only, title={#2}, minted language={#1}, minted style=xcode, minted options={fontsize=\fontsize{8}{9.6}, fontfamily=tt}, left=3mm}

\title{Python入门（六）}
\author{张博涵}
\institute{北京航空航天大学}
\date{\today}

\begin{document}

\begin{frame}[plain]
\maketitle
\end{frame}


\begin{frame}
    \frametitle{Contents}

    \tableofcontents

\end{frame}

\section{I/O}
\begin{frame}
    \frametitle{I/O Stream}

    在计算机系统中，不同程序之间、程序与系统之间以数据流(Stream)的形式传输数据，也可以称为 I/O Stream，I 代表 Input，O代表Output。

    \vspace{5mm}

    流一般有如下几种途径：

    \begin{outline}
        \1 标准流：标准输入(stdin)、标准输出(stdout)、标准错误(stderr)
        \2 在默认情况下，标准输入流为键盘，标准输出流和错误流为显示器（Shell），Python中通过\mintinline{python}{print}函数打印的内容传入标准输出，\mintinline{python}{input}函数会读取标准输入，报错信息传入标准错误流。
        \2 在Python中，Shell中均可以通过重定向(redirection)将输入流、输出流、错误流定向到文件
        \1 文件流：读取文件内容、向文件写入
    \end{outline}

    \vspace{5mm}

    流具有如下两种类型：

    \begin{itemize}
        \item 文本流：内容是字符串，Python会自动进行编码、解码、转义等，文本流本身有不同的编码方式（字符串章节提过）如UTF-8、GBK、GB2312。
        \item 二进制流：内容是二进制代码，不进行编码解码。
    \end{itemize}
\end{frame}

\begin{frame}[fragile]
    \frametitle{文件读取/写入}

    在Python中进行文件的读取和写入需要首先和文件建立连接，创建一个指向目标的I/O Stream，然后执行读/写操作。

    \begin{block}{读写文件}
        \begin{itemize}
            \item 可以使用\mintinline{python}{open}函数来连接某个文件，主要用法
            
            \begin{codebox}{Python}{open}
open(file_path, mode='r', encoding=None)\end{codebox}

            \item 这里仅列出部分常用参数，其他参数可以查找文档
            \item \mintinline{python}{file_path} 为文件路径
            \item \mintinline{python}{mode}为读写模式，可选包括\mintinline{python}{r, w, a, rb, wb, ab}六种
            \item \mintinline{python}{b}代表二进制流，在读写图像、视频等非文本类数据时常用，\mintinline{Python}{r, w, a}分别代表read(读取), write(写入), append(追加)，write和append的区别在于前者会删除文件现有的内容再写入，而后者在文件的末尾写入。
            \item \mintinline{python}{encoding}代表当流为文本流时的编码方式，默认为UTF-8。
        \end{itemize}
    \end{block}
\end{frame}

\begin{frame}[fragile]
    \frametitle{文件读取/写入}
\begin{codebox}{Python}{示例代码}
f = open('test.txt', 'w', encoding='UTF-8')
f.write('你好, test.txt!\n')
f.write('Hello again, test.txt!\n')
f.close()

f = open('test.txt', 'r', encoding='UTF-8')
contents = f.read()
print(contents)
f.close()
\end{codebox}

\begin{itemize}
    \item \mintinline{python}{open}函数返回一个文件对象，可以看做和文件的I/O接口
    \item 在写入模式下，可以通过\mintinline{python}{f.write()}像文件中写入数据
    \item 在读取模式下，可以通过\mintinline{python}{f.read()}读取全部文件内容，也可以通过\mintinline{python}{f.readlines(), f.readline()}等方法分行读取或逐行读取。
    \item 在流使用完毕后，应该借助\mintinline{python}{f.close()}方法关闭该连接，以免错误写入。
\end{itemize}

\end{frame}


\begin{frame}[fragile]
\frametitle{文件读取/写入}
\begin{codebox}{Python}{示例代码}
with open("test.txt", 'r') as f:
    contents = f.read()
print(contents)
\end{codebox}

\begin{itemize}
    \item 为了方便，避免忘记关闭连接，Python提供了\mintinline{python}{with ... as ...} 特殊语法。
    \item 可以在\mintinline{python}{with ... as ...}代码块内执行读取或写入的操作，当离开该代码块时，连接会被自动关闭。
    \item 不止是在文件读写方面，该代码块也在其他地方被使用。
\end{itemize}

\end{frame}


\begin{frame}
    \frametitle{字符串格式化}

    在Python中字符串的格式化方法主要有：

    \begin{itemize}
        \item f-string
        \item \mintinline{python}{'str{},{}'.format()}方法
        \item 其他常用函数
    \end{itemize}

\vspace{5mm}

    参考链接：

    \url{https://docs.python.org/3/tutorial/inputoutput.html}

\end{frame}


\section{错误处理: error handling}


\begin{frame}
    \frametitle{错误: Error}

    由于编程问题导致出现的程序错误，也通常被称为bug，错误主要有如下几种：

    \begin{itemize}
        \item 语法错误 (Syntax error)，这种错误最容易发现，因为语法错误时程序无法运行
        \item 逻辑错误 (logic error)，这种错误非常常见也是很容易犯的，可能会导致意料之外的结果但不一定会报出错误并导致程序终止，例如在设置条件时小于等于5设置为小于5
        \item 运行时错误 (runtime error)，这种错误也非常常见，例如访问只有5个元素的数组的第6个元素，或需要字符串时传入数字，这种错误可以通过测试来避免
    \end{itemize}

    错误处理的主要目的即是在合适的地方报出(raise)错误，并决定处理这种错误，如执行特定的代码，还是终止程序并报出错误信息以提醒用户，主要解决的错误类型为 runtime error。
   
    在Python中，error也常被称为Exception。

\end{frame}

\begin{frame}[fragile]
    \frametitle{Python中的Error}

\begin{codebox}{Python}{示例错误}
>>> 10 * (1/0)
Traceback (most recent call last):
    File "<stdin>", line 1, in <module>
ZeroDivisionError: division by zero
>>> 4 + spam*3
Traceback (most recent call last):
    File "<stdin>", line 1, in <module>
NameError: name 'spam' is not defined
>>> '2' + 2
Traceback (most recent call last):
    File "<stdin>", line 1, in <module>
TypeError: can only concatenate str (not "int") to str
\end{codebox}

\begin{itemize}
    \item \mintinline{python}{traceback}的主要目的是逐层从外到里给出错误信息的报出的地方
    \item \mintinline{python}{NameError, TypeError}等为Error的类型，在Python中，Error主要通过类和对象来实现，也可以通过继承基础的错误类型来实现自己的Error类型。
    \item 冒号后面的内容为错误信息，主要用于提示用户
\end{itemize}

\end{frame}


\begin{frame}[fragile]
    \frametitle{raise error: 报出错误}

    \begin{itemize}
        \item 在Python中最基础的错误类型为\mintinline{python}{BaseException}和\mintinline{python}{Exception}，其他的错误类型如\mintinline{Python}{ValueError}, \mintinline{python}{SyntaxError}等继承了这两类实现的。
        \item 报出错误首先需要一个错误类型的对象，然后使用\mintinline{python}{raise}函数报出，在不做任何处理的情况下，会中断程序。
        \item 在初始化对象时，第一个参数为报错信息。
    \end{itemize}

\begin{codebox}{Python}{示例代码}
raise(ValueError("this is a ValueError."))
raise(Exception("this is a normal Exception."))
\end{codebox}
\end{frame}


\begin{frame}[fragile]
\frametitle{处理错误}

可以通过\mintinline{python}{try ... except ...}语法来捕获(capture)并处理错误。

\begin{codebox}{Python}{示例代码}
while True:
    try:
        x = int(input("Please enter a number: "))
        break
    except xxx1:
        # some code
    except (xxx2, xxx3):
        # some code
    finally:
        # some code
    else:
        # some code\end{codebox}
\begin{outline}
    \1 当遇到\mintinline{python}{try}时，会首先执行\mintinline{python}{try}后面紧跟的代码块
    \1 若有错误报出，则根据错误的类型进行匹配并执行对应代码块中的代码，例如若错误类型为\mintinline{python}{ValueError}，则执行\mintinline{Python}{print}语句。
    \2 可以有多个\mintinline{python}{except ExceptionType:}分句。
    \2 若匹配不到，则执行\mintinline{python}{finally}语句(可选)，否则报出错误，终止程序
    \1 若没有错误，则执行\mintinline{python}{else}后面的代码(如果有的话)。
\end{outline}


\end{frame}





\end{document}